\documentclass[11pt,a4paper]{article}

% Template visual Matriz Aprovação. Compile com XeLaTeX ou LuaLaTeX.
\usepackage{fontspec}
\usepackage{polyglossia}
\setmainlanguage{portuguese}
\IfFontExistsTF{Aptos}{
  \setmainfont{Aptos}
  \setsansfont{Aptos}
}{
  \IfFontExistsTF{Calibri}{
    \setmainfont{Calibri}
    \setsansfont{Calibri}
  }{
    \IfFontExistsTF{TeX Gyre Heros}{
      \setmainfont{TeX Gyre Heros}
      \setsansfont{TeX Gyre Heros}
    }{
      \setmainfont{Latin Modern Sans}
      \setsansfont{Latin Modern Sans}
    }
  }
}
\usepackage[a4paper,margin=2cm,headheight=16pt]{geometry}
\usepackage{graphicx}
\usepackage{xcolor}
\usepackage{tikz}
\usepackage{tcolorbox}
\usepackage{enumitem}
\usepackage{booktabs}
\usepackage{tabularx}
\usepackage{amssymb}
\usepackage{fancyhdr}
\usepackage{titlesec}
\usepackage{hyperref}

\definecolor{matrizlime}{HTML}{CBFF4D}
\definecolor{matrizink}{HTML}{101313}
\definecolor{matrizmuted}{HTML}{596363}
\definecolor{matrizline}{HTML}{DDE3DF}
\definecolor{matrizsoft}{HTML}{F3F7EC}

\hypersetup{colorlinks=true,linkcolor=matrizink,urlcolor=matrizink}
\pagestyle{fancy}
\fancyhf{}
\lhead{\textsf{\small MATRIZ APROVAÇÃO}}
\rhead{\textsf{\small APOSTILA}}
\cfoot{\textsf{\small \thepage}}
\renewcommand{\headrulewidth}{0.4pt}
\renewcommand{\headrule}{\hbox to\headwidth{\color{matrizline}\leaders\hrule height \headrulewidth\hfill}}

\titleformat{\section}{\Large\bfseries\color{matrizink}}{\thesection}{0.6em}{}
\titleformat{\subsection}{\large\bfseries\color{matrizink}}{\thesubsection}{0.6em}{}
\setlist{nosep,leftmargin=*}

\newtcolorbox{foco}{colback=matrizsoft,colframe=matrizlime,boxrule=1pt,arc=3pt,left=10pt,right=10pt,top=8pt,bottom=8pt}
\newtcolorbox{lei}{colback=matrizink,colframe=matrizink,coltext=white,boxrule=0pt,arc=3pt,left=10pt,right=10pt,top=8pt,bottom=8pt}
\newtcolorbox{alerta}{colback=white,colframe=matrizline,boxrule=0.6pt,arc=3pt,left=10pt,right=10pt,top=8pt,bottom=8pt}

% Logo usada na capa. Caminho relativo à pasta onde o .tex é compilado
% (materiais/<disciplina>/). Ajuste se a estrutura mudar.
\newcommand{\logomatriz}{../../public/logo.png}

\begin{document}

\begin{titlepage}
  \thispagestyle{empty}
  \begin{center}
    \includegraphics[width=0.55\linewidth]{\logomatriz}\par
  \end{center}
  \vspace{2.2cm}
  {\sffamily\fontsize{28}{32}\selectfont\bfseries Apostila — Química: Fundamentos e estequiometria\par}
  \vspace{0.4cm}
  {\sffamily\large ESA, EsPCEx, AFA, ITA, IME, EFOMM, Escola Naval e ENEM\par}
  \vfill
  \begin{foco}
    \textsf{\textbf{Foco da apostila}}\\[3pt]
    balanceie a equação antes de fazer qualquer proporção. Os coeficientes indicam relações em mol; os índices da fórmula não podem ser alterados durante o balanceamento.
  \end{foco}
  \vspace{0.7cm}
  {\sffamily\small Atualizado em 16 de agosto de 2026\quad ·\quad adaptação necessária ao edital e ao nível da prova\par}
  \vspace{1cm}
  {\sffamily\small matrizaprova.com\par}
\end{titlepage}

\tableofcontents
\newpage

% Conteúdo gerado entra aqui.
\section*{O que cai}

As questões introdutórias de Química avaliam a leitura de fórmulas, a classificação de substâncias, a conservação da matéria e os cálculos estequiométricos. Mesmo quando a questão parece conceitual, a proporção entre quantidades costuma ser o passo decisivo.

Os pontos de maior incidência são:

\begin{enumerate}
  \item substâncias simples, compostas, misturas e métodos de separação;
  \item estrutura atômica, número atômico, número de massa e íons;
  \item distribuição eletrônica e formação de ligações;
  \item balanceamento e interpretação de equações químicas;
  \item mol, massa molar, reagente limitante e rendimento.
\end{enumerate}

\textbf{Regra de prova:} balanceie a equação antes de fazer qualquer proporção. Os coeficientes indicam relações em mol; os índices da fórmula não podem ser alterados durante o balanceamento.

\section*{Teoria essencial}

\subsection*{1. Matéria, substâncias e misturas}

Matéria é tudo que possui massa e ocupa lugar no espaço. Uma substância pura possui composição definida. Pode ser simples, formada por um único elemento químico, como O₂, ou composta, formada por mais de um elemento, como H₂O.

Mistura reúne duas ou mais substâncias. Mistura homogênea apresenta uma única fase visível, como água com sal dissolvido. Mistura heterogênea apresenta duas ou mais fases, como água e óleo. Fase é cada porção visualmente uniforme do sistema.

Não confunda fase com componente. Água e gelo formam um sistema heterogêneo com um único componente químico, a água, mas duas fases físicas.

Métodos de separação usam propriedades físicas:

\begin{tabularx}{\linewidth}{lX}
\toprule
 \textbf{Método} & \textbf{Aplicação típica} \\
\midrule
 filtração & sólido insolúvel + líquido \\
 decantação & fases com densidades diferentes \\
 centrifugação & acelera a sedimentação \\
 destilação simples & sólido dissolvido ou líquidos com grande diferença de ebulição \\
 destilação fracionada & líquidos miscíveis com pontos de ebulição próximos \\
 separação magnética & sólido atraído por ímã + outro sólido \\
 cromatografia & componentes com diferentes interações com o meio \\
\bottomrule
\end{tabularx}


Mudança física não altera a identidade química da substância. Fusão, ebulição e dissolução são exemplos usuais. Na mudança química, novas substâncias são formadas, como na combustão, na ferrugem e na decomposição.

\subsection*{2. Estrutura do átomo}

O núcleo contém prótons e nêutrons; a eletrosfera contém elétrons. O número atômico, representado por Z, é o número de prótons e identifica o elemento. O número de massa, representado por A, é a soma de prótons e nêutrons:

número de massa = prótons + nêutrons.

Em um átomo neutro, o número de elétrons é igual ao número de prótons. Íon é a espécie eletricamente carregada. Cátion resulta da perda de elétrons e possui carga positiva. Ânion resulta do ganho de elétrons e possui carga negativa.

Isótopos têm o mesmo número de prótons e diferentes números de nêutrons. Como pertencem ao mesmo elemento, possuem o mesmo Z, mas A diferente. Isóbaros têm o mesmo A e diferentes Z. Isótonos têm o mesmo número de nêutrons.

\subsection*{3. Tabela periódica e propriedades}

Elementos de uma mesma família apresentam propriedades químicas relacionadas, pois possuem configurações semelhantes na camada de valência. Os gases nobres formam a família 18 e apresentam baixa reatividade em condições usuais.

Em geral, o raio atômico aumenta para baixo em um grupo e para a esquerda em um período. A energia de ionização e a eletronegatividade tendem a aumentar para cima e para a direita, com exceções que podem ser indicadas pelo enunciado.

Metais tendem a perder elétrons e formar cátions. Não metais tendem a ganhar ou compartilhar elétrons. A eletronegatividade mede a tendência de atrair elétrons em uma ligação, com o flúor entre os elementos mais eletronegativos.

\subsection*{4. Ligações químicas}

Na ligação iônica, há transferência de elétrons, geralmente de um metal para um não metal. O resultado é uma rede de íons. Compostos iônicos costumam apresentar altos pontos de fusão e conduzem corrente quando fundidos ou dissolvidos em água, pois os íons podem se movimentar.

Na ligação covalente, os átomos compartilham pares de elétrons. Ela ocorre principalmente entre não metais. A molécula pode ser polar ou apolar, conforme a distribuição de cargas e a geometria. Uma molécula com ligações polares não é necessariamente polar: a simetria pode cancelar os dipolos.

Na ligação metálica, os cátions ficam organizados em uma estrutura com elétrons deslocalizados. Esse modelo explica a condução elétrica, a condução térmica e a maleabilidade dos metais.

Para obter a fórmula de um composto iônico, as cargas totais devem se compensar. Assim, Ca²⁺ e Cl⁻ formam CaCl₂; Al³⁺ e O²⁻ formam Al₂O₃. Os índices representam quantidades na fórmula e não devem ser trocados sem verificar a neutralidade.

\subsection*{5. Reações e equações químicas}

Uma equação química representa reagentes à esquerda e produtos à direita. A seta indica a transformação. Estados físicos podem ser indicados por (s), (l), (g) e (aq).

Balancear é ajustar coeficientes para que o número de átomos de cada elemento seja igual nos dois lados. A conservação da matéria exige essa igualdade.

Exemplo: a formação da água pode ser escrita como H₂ + O₂ → H₂O. O oxigênio tem dois átomos à esquerda e um à direita. A forma balanceada é 2 H₂ + O₂ → 2 H₂O. Foram alterados coeficientes, não os índices das fórmulas.

Em reações de combustão de hidrocarbonetos, balanceie primeiro carbono, depois hidrogênio e, por último, oxigênio. Em reações iônicas, pode ser útil separar os íons que efetivamente participam e cancelar os espectadores, desde que o nível da questão permita essa abordagem.

\subsection*{6. Mol e massa molar}

Mol é uma unidade de quantidade de matéria. Um mol contém aproximadamente 6,02 × 10²³ entidades elementares. A entidade pode ser átomo, molécula, íon ou unidade fórmula, conforme o caso.

A massa molar é a massa de um mol da espécie, expressa em g/mol. Para obtê-la, some as massas atômicas de todos os átomos presentes na fórmula, multiplicando cada massa pelo respectivo índice.

Exemplo: a massa molar da água é 2 × 1 + 16 = 18 g/mol. Portanto, 36 g de água correspondem a 36 ÷ 18 = 2 mol.

As relações fundamentais são:

\begin{itemize}
  \item quantidade de matéria = massa ÷ massa molar;
  \item número de entidades = quantidade de matéria × constante de Avogadro;
  \item massa = quantidade de matéria × massa molar.
\end{itemize}

Em gases, se o enunciado informar condições específicas, use o volume molar indicado. Não adote automaticamente 22,4 L/mol sem conferir temperatura e pressão, pois esse valor corresponde a condições de referência específicas.

\subsection*{7. Estequiometria}

O cálculo estequiométrico relaciona as quantidades de reagentes e produtos pela equação balanceada. O procedimento geral é:

\begin{enumerate}
  \item escrever e balancear a reação;
  \item converter o dado para mol, se ele estiver em massa ou volume;
  \item usar a proporção entre os coeficientes;
  \item converter o resultado para a unidade pedida.
\end{enumerate}

Exemplo: na reação 2 H₂ + O₂ → 2 H₂O, 2 mol de H₂ produzem 2 mol de H₂O. A proporção é 1 para 1 entre essas duas espécies, mas não entre H₂ e O₂: são necessários 2 mol de H₂ para cada 1 mol de O₂.

\subsection*{8. Reagente limitante e rendimento}

Quando dois ou mais reagentes são fornecidos, pode faltar um deles antes que o outro seja consumido. O reagente limitante determina a quantidade máxima de produto. O excesso permanece parcialmente sem reagir.

Para identificar o limitante, calcule quanto produto cada reagente poderia formar ou compare a quantidade disponível com a proporção exigida na equação. Use o limitante nos cálculos finais.

O rendimento compara a quantidade obtida na prática com a quantidade teórica:

rendimento percentual = quantidade real ÷ quantidade teórica × 100.

Se o rendimento for 80\%, a quantidade real será 80\% da quantidade teórica. Não aplique o percentual ao reagente antes de determinar o produto teórico quando o enunciado apresentar um limitante.

\section*{Como a banca cobra}

\begin{itemize}
  \item Confunde mistura homogênea com substância pura.
  \item Troca número de massa por número atômico.
  \item Trata cátion como espécie que ganhou elétrons.
  \item Permite alterar índice de fórmula durante o balanceamento.
  \item Usa coeficientes de uma equação não balanceada para fazer proporção.
  \item Esquece que a massa molar deve ser usada para converter gramas em mol.
  \item Escolhe o reagente em maior massa como limitante, sem comparar os mols e a proporção da reação.
\end{itemize}

\subsection*{Exemplos resolvidos}

\subsubsection*{Exemplo 1}

Determine a massa de água formada na reação entre 4 g de hidrogênio e oxigênio em excesso: 2 H₂ + O₂ → 2 H₂O. Use H = 1 e O = 16.

\begin{enumerate}
  \item Massa molar de H₂: 2 g/mol.
  \item Quantidade de H₂: 4 ÷ 2 = 2 mol.
  \item A equação indica 2 mol de H₂ para 2 mol de H₂O, portanto formam-se 2 mol de
\end{enumerate}

   água.

\begin{enumerate}
  \item Massa molar de H₂O: 18 g/mol.
  \item Massa formada: 2 × 18 = 36 g.
\end{enumerate}

\subsubsection*{Exemplo 2}

Na reação 2 CO + O₂ → 2 CO₂, estão disponíveis 3 mol de CO e 2 mol de O₂. Determine o reagente limitante.

\begin{enumerate}
  \item A proporção exige 2 mol de CO para 1 mol de O₂.
  \item Para consumir 3 mol de CO, seriam necessários 1,5 mol de O₂.
  \item Há 2 mol de O₂, quantidade suficiente para consumir todo o CO.
  \item Portanto, CO é o limitante e O₂ está em excesso.
\end{enumerate}

\section*{Quadro-resumo}

\begin{tabularx}{\linewidth}{llX}
\toprule
 \textbf{Conceito} & \textbf{Regra rápida} & \textbf{Erro comum} \\
\midrule
 número atômico & prótons & confundir com nêutrons \\
 número de massa & prótons + nêutrons & usar elétrons na soma \\
 cátion & perde elétrons & dizer que ganhou elétrons \\
 balanceamento & altera coeficientes & alterar índices \\
 mol & 6,02 × 10²³ entidades & confundir mol com grama \\
 massa molar & soma das massas da fórmula & esquecer índices \\
 limitante & acaba primeiro na proporção & escolher pela maior massa \\
 rendimento & real ÷ teórico × 100 & usar denominador errado \\
\bottomrule
\end{tabularx}


\section*{Questões de fixação}

\subsection*{Questão 1 — (questão inédita)}

Uma amostra contém água e etanol completamente misturados. O sistema é:

A) substância simples. B) substância composta. C) mistura homogênea. D) mistura heterogênea. E) elemento químico.

\begin{alerta}
\textbf{Gabarito: C.} Há mais de uma substância, mas apenas uma fase visível. Trata-se de uma mistura homogênea.
\end{alerta}

\subsection*{Questão 2 — (questão inédita)}

Um átomo neutro possui número atômico 17 e número de massa 35. O número de nêutrons e o número de elétrons são, respectivamente:

A) 17 e 18. B) 18 e 17. C) 35 e 17. D) 52 e 17. E) 17 e 35.

\begin{alerta}
\textbf{Gabarito: B.} Nêutrons = A − Z = 35 − 17 = 18. Como o átomo é neutro, o número de elétrons é igual ao número de prótons: 17.
\end{alerta}

\subsection*{Questão 3 — (questão inédita)}

Considere a reação não balanceada H₂ + O₂ → H₂O. Os menores coeficientes inteiros que a balanceiam são:

A) 1, 1, 1. B) 1, 2, 1. C) 2, 1, 2. D) 2, 2, 1. E) 1, 2, 2.

\begin{alerta}
\textbf{Gabarito: C.} Dois H₂ fornecem quatro átomos de hidrogênio e um O₂ fornece dois átomos de oxigênio. Dois H₂O contêm quatro hidrogênios e dois oxigênios.
\end{alerta}

\subsection*{Questão 4 — (questão inédita)}

Qual é a quantidade de matéria em 44 g de CO₂? Considere C = 12 e O = 16.

A) 0,5 mol. B) 1 mol. C) 2 mol. D) 22 mol. E) 44 mol.

\begin{alerta}
\textbf{Gabarito: B.} A massa molar do CO₂ é 12 + 2 × 16 = 44 g/mol. Logo, 44 g correspondem a 1 mol.
\end{alerta}

\subsection*{Questão 5 — (questão inédita)}

Na reação 2 H₂ + O₂ → 2 H₂O, uma mistura contém 3 mol de H₂ e 1 mol de O₂. Qual é o reagente limitante?

A) H₂, pois possui menor massa molar. B) H₂, pois são necessários 2 mol para cada 1 mol de O₂. C) O₂, pois possui maior massa molar. D) O₂, pois há apenas 1 mol disponível. E) Nenhum, pois os dois são completamente consumidos.

\begin{alerta}
\textbf{Gabarito: B.} Para 1 mol de O₂, a reação exige 2 mol de H₂. Como há 3 mol de H₂, sobra hidrogênio; o oxigênio é consumido primeiro e é o limitante.
\end{alerta}

\section*{Revisão final}

\begin{itemize}
  \item[$\square$] Diferenciei substância pura de mistura.
  \item[$\square$] Identifiquei Z, A, prótons, nêutrons e elétrons.
  \item[$\square$] Conferi a carga de cátions e ânions.
  \item[$\square$] Balanceei a equação sem alterar índices.
  \item[$\square$] Calculei a massa molar considerando todos os índices.
  \item[$\square$] Converti massa para mol antes da proporção.
  \item[$\square$] Verifiquei o reagente limitante pela equação balanceada.
  \item[$\square$] Separei quantidade teórica de quantidade real.
  \item[$\square$] Apliquei corretamente o percentual de rendimento.
\end{itemize}

\end{document}
